\documentclass[10pt,a4paper]{amsart}
\usepackage[margin=19mm]{geometry}
\usepackage{amsmath,amssymb,mathtools}
\usepackage[scr=rsfs,cal=euler]{mathalfa}
\usepackage{xfrac}
\usepackage[unicode,hidelinks,bookmarks=false]{hyperref}
\newcommand{\ie}{i.e.\ }
\newcommand{\cf}{cf.\ }
\newcommand{\resp}{respectively\ }
\newcommand{\Subsection}{\subsection}
\providecommand{\nobreakdash}{\nobreak\hbox{-}}
\setlength{\emergencystretch}{2em}
\multlinegap0pt
\usepackage{amssymb}
\newtheorem{theorem}{Theorem}
\newtheorem{lemma}{Lemma}
\newcommand{\LCA}{\mathrm{LCA}}
\newcommand{\MT}{\mathrm{MT}}
\DeclareMathOperator*{\argmin}{arg\,min}
\renewcommand{\epsilon}{\varepsilon}
\newcommand{\im}{\mathrm{im}\,}
\title{Intrinsic Bottleneck Distance for Merge Trees}
\author{David Beers\; and Gillian Grindstaff}
\date{Source: arXiv:2509.02755v2 (CC BY 4.0)}
\begin{document}
\maketitle

\noindent\emph{Source and licence.} Intrinsic Bottleneck Distance for Merge Trees. Version arXiv:2509.02755v2 (CC BY 4.0), distributed by arXiv under the Creative Commons Attribution 4.0 International licence, http://creativecommons.org/licenses/by/4.0/, as recorded in the arXiv licence metadata for this exact version and shown on the abstract page https://arxiv.org/abs/2509.02755v2. Full licence notice: \texttt{licenses/arxiv-2509.02755-CC-BY-4.0.txt}.\par\smallskip

\noindent\textit{Editorial scope.} Counted unit: the article's theorem thm:localeq (source0.tex lines 323--329). The article's complete original proof of it is delivered with it (source0.tex lines 629--710, one \texttt{proof} environment of 82 lines, which the article places away from the statement). No mathematics is written, repaired or re-derived: the statement and the proof are the article's own lines, in the article's own notation, and no retained line is substituted or re-flowed.  Retained uncounted as the source-local context the proof needs: lines 235--238; lines 538--549.  Nothing was omitted from any retained span, so the package carries no omission record.  No retained line contains a citation, so the wrapper carries no bibliography and no bibliography entry.  No retained line contains a figure environment, an image-inclusion command or drawing code, so no image is delivered and the package declares no asset dependency.  Editorial formatting only, in addition: an AMS \texttt{amsart} preamble that restates the article's own declarations and macros used by the retained lines, namely the environments \texttt{theorem}, \texttt{lemma} on separate counters, together with the standard packages those lines need.  The retained environments print in this document as Theorem 1, Lemma 1 in order of appearance, whereas the article prints the same items with the numbering of its own full document.  The complete native source of the article as distributed in the arXiv e-print for this exact version is bundled unchanged as \texttt{sources/184356/source0.tex}.
\begin{theorem}
    \label{thm:bottlelessthaninterleaving}
    On the space $\MT$, $d_B \leq d_I$.
\end{theorem}

\begin{lemma}
    \label{lem:plemma}
    For each $j \in E$ there exists a unique solution to the optimization problem
    \begin{equation*}
        \argmin_{x \in T'} f'(x) \quad \text{subject to} \quad l'_{\sigma(i)}, l'_j \preceq x \text{, for some }\sigma(i) < j \,.
    \end{equation*}
    Denote this solution by $p'_j$. We have that each $p'_j$ is a branch point and satisfies
    \begin{equation*}
        f'(p'_j) - f'(l'_j) \leq 2\delta\,.
    \end{equation*}
    Let $q$ be the least index such that $l'_q \preceq p'_j$. Then $q \in \im \sigma$.
\end{lemma}

\begin{theorem}
    \label{thm:localeq}
    Let $(T,f)\in\MT$. There is a neighborhood $W$ of $(T,f)$ in $\MT$ such that for all $(T',f')\in W$,
    \begin{equation*}
        d_B\big((T,f),(T',f')\big) = d_I\big((T,f),(T',f')\big)\,.
    \end{equation*}
\end{theorem}

\begin{proof}[Proof of Theorem \ref{thm:localeq}.] For $x = i^t(l_i)$, define
\begin{equation*}
    \widetilde{\alpha}_i(x) = i^{t + \delta +  f(l_i) - f'(l'_{\sigma(i)})}(l'_{\sigma(i)})\,.
\end{equation*}
This definition makes sense since $t + \delta +  f(l_i) - f'(l'_{\sigma(i)}) \geq t \geq 0$. Moreover, $f'\widetilde{\alpha}_i(x) = f(x) + \delta$.
If $x = i^t(l_i) = i^s(l_j)$, then $\LCA(l_i,l_j) \preceq x$. Therefore,
\begin{equation*}
    f'\widetilde{\alpha}_i(x) = f(x) + \delta \geq f(\LCA(l_i,l_j)) + \delta \geq f'(\LCA(l'_{\sigma(i)}, l'_{\sigma(j)}))\,.
\end{equation*}
Since $\LCA(l'_{\sigma(i)}, l'_{\sigma(j)})$ and $\widetilde{\alpha}_i(x)$ are both ancestors of $l'_{\sigma(i)}$, this means $\LCA(l'_{\sigma(i)}, l'_{\sigma(j)}) \preceq \widetilde{\alpha}_i(x)$. Similarly, $\LCA(l'_{\sigma(i)}, l'_{\sigma(j)}) \preceq \widetilde{\alpha}_j(x)$. Since $f'\widetilde{\alpha}_i(x) = f'\widetilde{\alpha}_j(x) = f(x) + \delta$, this implies $\widetilde{\alpha}_i(x) = \widetilde{\alpha}_j(x)$. Therefore the maps $\widetilde{\alpha}_i$ together define a map $\widetilde{\alpha}: T \to T'$. Since $f'\widetilde\alpha(x) = f(x) + \delta$ and $\widetilde\alpha(x) \preceq \widetilde\alpha(y)$ whenever $x \preceq y$, we deduce that $\widetilde\alpha$ is continuous.

Due to the last statement of Lemma \ref{lem:plemma}, for $k \in E$ we may define $w(k)$ such that $\sigma(w(k))$ is the least index of a leaf descendant from $p'_k$. With this in mind, for $x = i^t(l'_k)$, we define
\begin{equation*}
    \widetilde{\beta}_k(x) = \begin{cases}
        i^{t + \delta + f'(l'_k) - f(l_i)}(l_i)& k = \sigma(i)\,,\\
        i^{t + \delta + f'(l'_k) - f(l_{w(k)})}(l_{w(k)})& k \in E\,.
    \end{cases}
\end{equation*}

If $k = \sigma(i)$, this is well defined as $t + \delta + f'(l'_k) - f(l_i) \geq t \geq 0$, and moreover we have $f\widetilde{\beta}_k(x) = t + \delta + f'(l'_k) = f'(x) + \delta$.

If $k \in E$, $\widetilde\beta_k(x)$ is well defined as
\begin{align*}
    t + \delta + f'(l'_k) - f(l_{w(k)}) &= t + \delta + f'(l'_{\sigma(w(k))}) - f'(l'_{\sigma(w(k))}) + f'(l'_k) - f(l_{w(k)})\\
    &\geq t + \delta + f'(l'_{\sigma(w(k))}) - f(l_{w(k)})\\
    &\geq t \geq 0\,,
\end{align*}
where in going from the first to the second line we have used that since $l'_{\sigma(w(k))}, l'_k\preceq p'_k$, we have $f'(l'_{\sigma(w(k))}) \leq f'(l'_k)$, by the definition of $w$. We also have $f\widetilde\beta_k(x) = t + \delta + f'(l'_{k}) = f'(x) + \delta$.

We now wish to show that the maps $\widetilde\beta_k$ combine to define a map $\widetilde\beta:T' \to T$, as we did with $\widetilde\alpha$. Let $x = i^t(l'_k) = i^s(l'_q)$. There are three cases to handle.

The first case is that $k = \sigma(i)$ and $q = \sigma(j)$. In this case $\LCA(l'_{\sigma(i)}, l'_{\sigma(j)})\preceq x$, so
\begin{equation*}
    f\widetilde\beta_k(x) = f'(x) + \delta \geq f'(\LCA(l'_{\sigma(i)}, l'_{\sigma(j)})) + \delta \geq f(\LCA(l_i,l_j))\,.
\end{equation*}
Since both $\widetilde\beta_k(x)$ and $\LCA(l_i,l_j)$ are ancestors of $l_i$, it follows that $\LCA(l_i,l_j) \preceq \widetilde \beta_k(x)$. Similarly, $\LCA(l_i,l_j) \preceq \widetilde \beta_q(x)$. Since $f\widetilde\beta_k(x) = f\widetilde\beta_q(x) = f'(x) + \delta$, we conclude that $\widetilde\beta_k(x) = \widetilde\beta_q(x)$.

The second case is that $k = \sigma(i)$ and $q \in E$. If $x \preceq p'_q$, then it must be the case that $f'(l'_{\sigma(i)}) > f'(l'_q)$. Therefore, using that $f'(p'_q) - f'(l'_q) \leq 2\delta \leq 2\epsilon$, we have
\begin{align*}
    f'(l'_q) < f'(l'_{\sigma(i)}) \leq f'(x) \leq f'(p'_q) \implies f'(x) - f'(l'_{\sigma(i)}) \leq 2\epsilon \implies f\beta(x) - f(l_i) < 4\epsilon\,,
\end{align*}
so that $l_i \preceq \beta(l'_{\sigma(i)}) \preceq \beta(x) \preceq i^{4\epsilon}(l_i)$. Since there are no branch points in $f^{-1}(f(l_i),f(l_i) + 4\epsilon]$, we have that all descendants of $\beta(x)$ are ancestors of $l_i$. In particular, $l_i \preceq \beta(l'_q)$. This implies $\beta(l'_{\sigma(i)}) \preceq \beta(l'_q)$, and so $f'(l'_q) \geq f'(l'_{\sigma(i)})$, a contradiction.

It follows that $p'_q \preceq x$, so that $x = i^u(l'_{\sigma(w(q))})$. We have that $u = s + f'(l'_q) - f'(l'_{\sigma(w(q))})$, so
\begin{align*}
    \widetilde\beta_{\sigma(w(q))}(x) &= i^{u + \delta + f'(l'_{\sigma(w(q))}) - f(l_{w(q)})}(l_{w(q)})\\
    & = i^{s + \delta + f'(l'_q) - f'(l'_{\sigma(w(q))}) + f'(l'_{\sigma(w(q))}) - f(l_{w(q)})}(l_{w(q)})\\
    &= i^{s + \delta + f'(l'_q) - f(l_{w(q)})}(l_{w(q)})\\
    & = \widetilde\beta_{q}(x)\,.
\end{align*}
Hence, by the first case, $\widetilde\beta_k(x) = \widetilde\beta_{\sigma(w(q))}(x) = \widetilde{\beta}_{q}(x)$.

The final case is that $k, q \in E$. In this case, if $l'_{\sigma(i)} \preceq x$ for some $1 \leq i \leq n$, then $x = i^u(l'_{\sigma(i)})$ and so, using the second case, we have $\widetilde\beta_{k}(x) = \widetilde\beta_{\sigma(i)}(x) = \widetilde\beta_{q}(x)$. Hence we may assume there is no $1 \leq i \leq n$ such that $l'_{\sigma(i)} \preceq x$. As such, $x \preceq p'_k$, so $l'_q \preceq x \preceq p'_k$\,. From this we deduce that $p'_k = p'_q$. Hence $w(k) = w(q)$. Using this and that $t + f'(l'_k) = s + f'(l'_q)$, we have
\begin{align*}
    \widetilde{\beta}_k(x) &= i^{t + \delta + f'(l'_k) - f(l_{w(k)})}(l_{w(k)})\\
    &= i^{s + \delta + f'(l'_q) - f(l_{w(q)})}(l_{w(q)})\\
    & = \widetilde\beta_q(x)\,.
\end{align*}

Hence the $\widetilde\beta_k$ together define a map $\widetilde\beta: T' \to T$. Since $f\widetilde\beta(x) = f'(x) + \delta$ and $\widetilde\beta(x) \preceq \widetilde\beta(y)$ whenever $x \preceq y$, $\widetilde\beta$ is continuous.

It only remains to verify the identities $\widetilde\alpha\widetilde\beta = i^{2\delta}$ and $\widetilde\beta\widetilde\alpha = i^{2\delta}$. If $x = i^t(l_i)$, then
\begin{align*}
    \widetilde\beta\widetilde\alpha(x) = \widetilde\beta i^{t + \delta +  f(l_i) - f'(l'_{\sigma(i)})}(l'_{\sigma(i)}) = i^{t + \delta +  f(l_i) - f'(l'_{\sigma(i)}) + \delta + f'(l'_{\sigma(i)}) - f(l_i)} (l_i) = i^{2\delta} i^t(l_i) = i^{2\delta}(x)\,.
\end{align*}

If $x \in T'$, first assume $x = i^t(l'_{\sigma(i)})$, then
\begin{align*}
    \widetilde\alpha\widetilde\beta(x) = \widetilde\alpha i^{t + \delta +  f'(l'_{\sigma(i)}) - f(l_{i})}(l_i) = i^{t + \delta +  f'(l'_{\sigma(i)}) - f(l_{i}) + \delta + f(l_i) - f'(l'_{\sigma(i)})} (l'_{\sigma(i)}) = i^{2\delta} i^t(l'_{\sigma(i)}) = i^{2\delta}(x)\,.
\end{align*}
Otherwise, $x = i^t(l'_j)$, for some $j \in E$. In this case we have $f'\widetilde\alpha\widetilde\beta(x) = t + 2\delta + f'(l'_j) = f'i^{2\delta}(x)$. Since $f'(p'_j) - f'(l'_j) \leq 2\delta$ and $l'_j$ is a descendant of $p'_j$, it follows that $l'_{\sigma(w(j))} \preceq p'_j \preceq i^{2\delta}(l'_j) \preceq i^{t+2\delta}(l'_j) = i^{2\delta}(x)$. Since $\widetilde\alpha\widetilde\beta(x)$ is also an ancestor of $l'_{\sigma(w(j))}$ with the same $f$-value as $i^{2\delta}(x)$, it follows that $\widetilde\alpha\widetilde\beta(x) = i^{2\delta}(x)$.

Hence $\widetilde\alpha$ and $\widetilde\beta$ define a $\delta$-interleaving between $(T,f)$ and $(T',f')$, establishing that
\begin{equation*}
    d_B\big((T,f),(T',f')\big) \geq d_I\big((T,f),(T',f')\big)\,.
\end{equation*}
Then Theorem \ref{thm:bottlelessthaninterleaving} implies that in fact
\begin{equation*}
    d_B\big((T,f),(T',f')\big) = d_I\big((T,f),(T',f')\big)\,,
\end{equation*}
and the proof is complete.
\end{proof}

\end{document}
